% Glossar-Einträge, alphabetisch sortiert.
% Format: \abk{Kürzel}{Langform}{Erklärung}
% Ohne Erklärung: \abk{Kürzel}{Langform}{}

\abk{ASIC}{Application-Specific Integrated Circuit}
  {Ist ein IC, der speziell für eine klar definierte Aufgabe entwickelt wurde. Er ist nicht frei programmierbar.}

\abk{ASSP}{Application-Specific Standard Product}
  {Ist ein IC, der für einen bestimmten Anwendungsbereich entwickelt wurde, im Gegensatz zu einem ASIC aber frei verkäuflich ist.}

\abk{FPGA}{Field-Programmable Gate Array}
  {Integrierter Schaltkreis, dessen Logikfunktion erst nach der Fertigung durch eine Konfigurationsdatei festgelegt wird.}

\abk{IC}{Integrated Circuit}
  {Eine elektronische Schaltung, die auf einem kleinen Plättchen aus Halbleitermaterial (meist Silizium, engl. Silicon) aufgebracht ist.}

\abk{ISA}{Instruction Set Architecture}{
Befehlssatzarchitektur. Legt fest, welche Befehle, Datentypen und Adressierungsarten ein Prozessor versteht. Beispiele: x86, x86-64, ARM.}

\abk{MIPS}{Million Instructions Per Second}
  {Maßeinheit für die Rechenleistung eines Prozessors.}

\abk{NRE}{Non-Recurring Engineering}
  {Einmalige Kosten für Entwicklung, Verifikation und Fertigungsvorbereitung, die unabhängig von der Stückzahl anfallen.}

\abk{TTM}{Time-to-Market}
  {Zeitspanne von der Produktidee bis zur Markteinführung.}
